% cap5_rezultate.tex
\chapter{Rezultate și evaluare}
\label{ch:rezultate}

Prezentul capitol descrie metodologia de evaluare adoptată pentru prototipul de sistem smart home
local-first și prezintă rezultatele obținute în faza curentă de dezvoltare. Sistemul se află la
stadiul de prototip funcțional asamblat pe breadboard; modelul de \gls{ml} bazat pe
\gls{if} nu a fost încă antrenat pe date reale de lungă durată, iar suita de teste automate
este în curs de implementare. Prin urmare, orice cifră numerică prezentată în acest capitol este
marcată explicit ca \textbf{valoare preliminară / ilustrativă} și reflectă fie măsurători punctuale
obținute în sesiuni de testare manuală, fie estimări bazate pe specificațiile tehnice ale
componentelor utilizate. Acestea urmează să fie validate statistic pe o perioadă mai lungă de
operare.

% ──────────────────────────────────────────────────────────────────────
\section{Mediul și metodologia de testare}
\label{sec:mediu-testare}
% ──────────────────────────────────────────────────────────────────────

\subsection{Configurația hardware de testare}

Întregul prototip a fost asamblat pe breadboard, fără circuit imprimat dedicat, ceea ce introduce
variabilitate suplimentară în măsurătorile de tensiune și timpi de răspuns. Componentele principale
sunt:

\begin{itemize}
  \item un nod interior bazat pe ESP32 WROOM-32, cu senzorii \gls{dht}11 (×2), MQ-2,
        \gls{ldr} (×2), buzzer, servomotor SG90 și trei \gls{led}-uri;
  \item un nod exterior ESP32-CAM (AI-Thinker, OV2640), cu senzor \gls{pir} pe GPIO13;
  \item un gateway Raspberry Pi 3B+ rulând Mosquitto (port 1883), PostgreSQL + TimescaleDB,
        \gls{api}-ul NestJS (port 3000) și motorul \gls{ml} Python.
\end{itemize}

Rețeaua locală de testare este un router Wi-Fi de uz casnic; descoperirea serviciilor se realizează
prin \gls{mdns} (\texttt{smarthome.local}, \texttt{smarthome-cam.local}).

\subsection{Modul simulator}

Deoarece senzorii fizici produc date cu rată de 5 secunde și necesită prezența fizică pentru
stimularea unor scenarii (de exemplu gaz detectat), \gls{api}-ul NestJS include un modul
\texttt{simulator/} activat prin variabila de mediu \texttt{SIMULATE=true}. Acesta generează
citiri sintetice cu variații realiste ale tuturor câmpurilor (inclusiv \texttt{temperature2},
\texttt{humidity2}, \texttt{light1}, \texttt{light2}), permițând testarea fluxurilor de date,
a persistenței în baza de date și a comunicării \gls{mqtt} fără hardware conectat.

\subsection{Piramida de testare propusă}

Metodologia de testare urmează structura clasică în trei niveluri, ilustrată în
\figref{fig:piramida}.

\begin{figure}[H]
  \centering
  \includegraphics[width=\singlefigure]{\pics/fig-piramida-testare}
  \caption{Piramida de testare propusă pentru sistem: teste unitare la bază, teste de integrare la
    nivel intermediar și teste end-to-end la vârf. Lateral — testarea manuală pe hardware.}
  \label{fig:piramida}
\end{figure}

\textbf{Nivelul 1 — Teste unitare.}
Serviciile \gls{api}-ului NestJS (autentificare, publicare comenzi, procesare alertă) și modulul
\texttt{features.py} (calculul celor 12 trăsături, normalizarea cu StandardScaler) sunt vizate
pentru testare cu Jest, respectiv \texttt{pytest}. Aceste teste verifică logica internă a fiecărei
componente în izolare, cu dependențe substituibile prin mock-uri.

\textbf{Nivelul 2 — Teste de integrare.}
Scenariile de integrare vizează lanțul complet: \gls{api} NestJS $\leftrightarrow$ broker
Mosquitto $\leftrightarrow$ baza de date PostgreSQL/TimescaleDB. Se verifică, de exemplu,
că o citire publicată pe topicul \texttt{smarthome/esp32\_node\_a/sensors} este persistată corect
în hypertable-ul \texttt{sensor\_readings} și propagată prin WebSocket pe canalul
\texttt{/live}.

\textbf{Nivelul 3 — Teste end-to-end.}
Scenariile end-to-end acoperă lanțul integral firmware $\leftrightarrow$ broker $\leftrightarrow$
\gls{api} $\leftrightarrow$ aplicație mobilă. Se verifică că o comandă emisă din interfața mobilă
ajunge la actuatorul fizic în sub-secundă și că o alertă de gaz generată de firmware apare pe
ecranul utilizatorului.

\textbf{Testare manuală pe hardware.}
Paralel cu testele automate (în curs de implementare), fiecare funcționalitate a fost verificată
manual: conectarea nodului la rețea, publicarea datelor \gls{mqtt}, executarea comenzilor de
actuator, generarea alertei de gaz și vizualizarea fluxului video de la camera ESP32-CAM.

% ──────────────────────────────────────────────────────────────────────
\section{Verificarea funcțională a sistemului}
\label{sec:verificare-functionala}
% ──────────────────────────────────────────────────────────────────────

Tabelul~\tabref{tab:verif-functionala} centralizează funcționalitățile principale ale sistemului,
metoda de verificare aplicată și statusul curent. Verificarea „manuală confirmat" înseamnă că
funcționalitatea a fost testată în sesiuni de lucru pe prototipul fizic sau cu simulatorul activat.

\begin{table}[H]
  \centering
  \caption{Verificarea funcțională a componentelor sistemului.}
  \label{tab:verif-functionala}
  \small
  \begin{tabular}{p{4.8cm} p{5.4cm} p{2.6cm}}
    \toprule
    \textbf{Funcționalitate} & \textbf{Metodă de verificare} & \textbf{Status} \\
    \midrule
    Conectare ESP32 la broker \gls{mqtt} (mDNS) &
      Monitorizare log Mosquitto; LED de stare GPIO25 &
      Confirmat manual \\
    \addlinespace
    Publicare citiri senzori la 5\,s pe topicul \texttt{sensors} &
      Subscriber MQTT extern (\texttt{mosquitto\_sub}); vizualizare dashboard &
      Confirmat manual \\
    \addlinespace
    Persistare în TimescaleDB (hypertable) &
      Interogare SQL pe \texttt{sensor\_readings}; endpointul \texttt{/api/sensors/history} &
      Confirmat manual \\
    \addlinespace
    Comanda \texttt{led\_on/off/toggle} (LED 0, 1, 2) &
      Observare vizuală; log firmware serial &
      Confirmat manual \\
    \addlinespace
    Comanda \texttt{fan\_on/off} (ventilator prin tranzistor) &
      Observare fizică + log serial &
      Confirmat manual \\
    \addlinespace
    Comanda \texttt{servo\_move} (0°–180°, draperii) &
      Observare mișcare servomotor SG90 &
      Confirmat manual \\
    \addlinespace
    Comanda \texttt{buzzer\_beep} / \texttt{all\_off} &
      Observare sonoră + log serial &
      Confirmat manual \\
    \addlinespace
    Alertă gaz (\texttt{gasLevel > 350}) → \texttt{smarthome/alerts} + beep &
      Injectare nivel gaz simulat; observare alertă în aplicație &
      Confirmat manual \\
    \addlinespace
    Propagare alertă prin WebSocket → aplicație mobilă &
      Vizualizare tab Alerte în timp real &
      Confirmat manual \\
    \addlinespace
    Stare nod (Last Will \gls{mqtt}, \texttt{status} retained) &
      Oprire bruscă ESP32; verificare flag \texttt{online=false} &
      Confirmat manual \\
    \addlinespace
    Stream video ESP32-CAM (\gls{mjpeg}) + snapshot &
      Acces \url{http://smarthome-cam.local/stream} din WebView &
      Confirmat manual \\
    \addlinespace
    Autentificare \gls{jwt} + bcrypt &
      Testare endpoint \texttt{/api/auth/login}; token valid/expirat &
      Confirmat manual \\
    \addlinespace
    Provisioning Wi-Fi (hotspot \texttt{SmartHome-Setup}) &
      Conectare la AP fallback; trimitere comandă \texttt{wifi\_update} &
      Confirmat manual \\
    \addlinespace
    Inferență \gls{ml} (cold-start, reguli prag) &
      Injectare valori extreme; verificare alertă \texttt{ML\_ANOMALY} &
      Confirmat manual \\
    \addlinespace
    Simulator \texttt{SIMULATE=true} — date mock &
      Pornire \gls{api} în mod simulator; verificare dashboard &
      Confirmat manual \\
    \addlinespace
    Teste automate (Jest / pytest) &
      Suită în curs de implementare &
      \textit{În lucru} \\
    \bottomrule
  \end{tabular}
\end{table}

\subsection{Fluxul de comandă}

Fluxul de comandă pornește de la interacțiunea utilizatorului în aplicația mobilă (de exemplu,
apăsarea butonului „LED 1 ON") și traversează: \gls{api} NestJS → publicare \gls{mqtt} pe topicul
\texttt{smarthome/esp32\_node\_a/commands} → recepție de către firmware → execuție \texttt{setLed(0,true)}.
Confirmarea nu este transmisă explicit înapoi (firmware-ul nu publică un ACK), dar starea
actuală se reflectă în citirea imediat următoare dacă un senzor asociat este prezent (de exemplu,
schimbarea iluminării detectată de \gls{ldr}).

\subsection{Fluxul de alertă de gaz}

Detecția gazului funcționează pe front crescător: firmware-ul evaluează la fiecare citire dacă
\texttt{gasLevel} (valoare \gls{adc} pe 10 biți, 0–1023) depășește pragul de 350. La depășire,
dacă alerta nu era deja activă (\texttt{gasAlert == false}), firmware-ul publică imediat pe
\texttt{smarthome/alerts}, activează buzzer-ul local (beep autonom, independent de rețea) și
setează \texttt{gasAlert = true}. Alerta dispare (front descrescător) când nivelul scade sub 350.
\gls{api}-ul persistă evenimentul ca \texttt{GAS\_DETECTED} și îl propagă prin WebSocket.

% ──────────────────────────────────────────────────────────────────────
\section{Performanță operațională}
\label{sec:performanta}
% ──────────────────────────────────────────────────────────────────────

\textit{Notă: valorile din această secțiune sunt \textbf{preliminare / ilustrative}, obținute
din sesiuni de testare manuală punctuală pe prototipul fizic sau estimate pe baza specificațiilor
tehnice. Ele nu au caracter statistic și urmează să fie validate prin măsurători sistematice pe
o perioadă extinsă.}

\subsection{Latența comandă–actuator}

Latența este definită ca intervalul dintre momentul în care aplicația mobilă emite cererea
\texttt{POST /api/commands/\{nodeId\}} și momentul în care firmware-ul execută acțiunea (observabil
prin \texttt{Serial.print} pe UART). Factorii care o influențează includ latența \gls{wifi} în
\gls{lan}, timpul de procesare al broker-ului Mosquitto și bucla principală a firmware-ului
Arduino.

\subsection{Frecvența de publicare senzori}

Intervalul de publicare este configurat în firmware la 5 secunde (\texttt{PUBLISH\_INTERVAL = 5000\,ms}).
Citirea senzorilor \gls{dht}11 durează tipic 250\,ms per senzor conform specificațiilor
producătorului~\cite{aosong2019dht11}, deci ciclul complet (2 × DHT11 + MQ-2 + 2 × \gls{ldr})
se încadrează confortabil sub pragul de 5 secunde.

\begin{table}[H]
  \centering
  \caption{Indicatori de performanță operațională — valori \textbf{preliminare / ilustrative}.}
  \label{tab:performanta}
  \small
  \begin{tabular}{p{5.0cm} p{3.2cm} p{4.4cm}}
    \toprule
    \textbf{Indicator} & \textbf{Valoare ilustrativă} & \textbf{Observații} \\
    \midrule
    Latență comandă→actuator (\gls{lan}) &
      $\approx 120$–$250\,\text{ms}$ &
      Măsurată manual (timestamp log serial); variabilă în funcție de încărcare Wi-Fi \\
    \addlinespace
    Interval publicare citiri senzori &
      $5\,\text{s}$ (configurat) &
      Parametru constant în firmware; citire DHT11 $\approx 250\,\text{ms}$/senzor \\
    \addlinespace
    Timp răspuns \gls{api} (\texttt{/api/sensors/latest}) &
      $< 50\,\text{ms}$ &
      Măsurare locală cu \texttt{curl}; fără latență de rețea externă \\
    \addlinespace
    Consum \gls{cpu} \gls{rpi} (idle, toate serviciile active) &
      $\approx 15$–$25\%$ &
      Estimare \texttt{htop}; vârfuri la reantrenare ML (03:00) \\
    \addlinespace
    Consum \gls{ram} \gls{rpi} (toate serviciile active) &
      $\approx 450$–$550\,\text{MB}$ &
      PostgreSQL + NestJS + Mosquitto + Python; din totalul de 1\,GB \\
    \addlinespace
    Timp inferență \gls{ml} per citire (\texttt{infer.py}) &
      $< 5\,\text{ms}$ &
      Estimare pe baza caracteristicilor \gls{if}; neglijabil față de intervalul de 5\,s \\
    \addlinespace
    Latență stream video (\gls{mjpeg}) în \gls{lan} &
      $\approx 300$–$800\,\text{ms}$ &
      Observare subiectivă în WebView; dependentă de rezoluție OV2640 \\
    \bottomrule
  \end{tabular}
\end{table}

Rezultatele preliminare indică faptul că sistemul este operabil pe hardware de uz general (\gls{rpi}
3B+) fără a atinge limitele de resurse în condiții normale de funcționare. Vârfurile de \gls{cpu}
la reantrenarea zilnică a modelului \gls{ml} (programată la ora 03:00 prin \texttt{scheduler.py})
nu afectează funcționarea în timp real datorită programării în afara intervalului de utilizare.

% ──────────────────────────────────────────────────────────────────────
\section{Evaluarea detecției de anomalii}
\label{sec:eval-anomalii}
% ──────────────────────────────────────────────────────────────────────

Detecția de anomalii comportamentale cu \gls{if} reprezintă contribuția originală a acestei
lucrări. Aceasta este prezentată detaliat din perspectivă algoritmică în \chapref{ch:metodologie}.
Prezenta secțiune descrie \textbf{protocolul de evaluare propus} și discută parametrii cheie ai
modelului, cu mențiunea că evaluarea cantitativă completă urmează a fi realizată după acumularea
unui set de date real de cel puțin 30 de zile de operare continuă.

\subsection{Protocolul de evaluare propus}

Evaluarea modelului \gls{if} se va realiza prin \textbf{injectare de anomalii sintetice} în date
istorice reale. Procedura propusă cuprinde pașii următori:

\begin{enumerate}
  \item \textbf{Colectarea datelor de referință.} Se colectează minim 30 de zile de citiri
    normale în condiții de viață reală (\texttt{sensor\_readings}), cu o rată de 5 secunde,
    reprezentând un set de referință de cel puțin $\sim$500.000 de rânduri.

  \item \textbf{Antrenarea modelului de referință.} Se antrenează modelul cu parametrii actuali:
    \texttt{IsolationForest(contamination=0.05, n\_estimators=100)} pe datele de referință,
    cu normalizare StandardScaler pe cele 12 trăsături din \texttt{features.py}.

  \item \textbf{Generarea anomaliilor sintetice.} Se injectează perturbații controlate în setul
    de test:
    \begin{itemize}
      \item spike-uri de temperatură ($+15\,°C$ față de medie, durată 1–5 minute),
      \item nivel de gaz ridicat artificial (\texttt{gasLevel > 600}, necorelat cu temperatura),
      \item mișcare detectată în interval de noapte (23:00–06:00) în absența altor activități,
      \item combinații atipice de trăsături (de exemplu umiditate 95\% corelată cu temperatură scăzută).
    \end{itemize}

  \item \textbf{Evaluarea cu metrici standard.} Se calculează precizia ($P$), recall-ul ($R$),
    scorul $F_1$ și curba \gls{roc} cu aria \gls{auc}-ROC~\cite{fawcett2006roc,davis2006prcurve}:
    \[
      F_1 = \frac{2 \cdot P \cdot R}{P + R}, \qquad
      \text{AUC-ROC} = \int_0^1 \text{TPR}(u)\,\mathrm{d}u.
    \]
    Alegerea pragului optimal se realizează prin analiza curbei Precizie-Recall~\cite{davis2006prcurve},
    preferată față de curba ROC în seturile dezechilibrate (anomaliile sunt rare).

  \item \textbf{Sensibilitate la prag.} Se variază \texttt{ANOMALY\_THRESHOLD} în intervalul
    $[-0.3, 0.0]$ pentru a cuantifica compromisul fals-pozitive vs.\ fals-negative.
\end{enumerate}

\subsection{Parametrii actuali ai modelului}

Modelul curent, descris în \texttt{gateway/ml/train.py}, folosește:
\begin{itemize}
  \item \textbf{contamination = 0.05}: se presupune că 5\% din datele de antrenament pot fi
    anomalii; aceasta determină pragul intern al arborelui de izolare~\cite{liu2008isolation};
  \item \textbf{ANOMALY\_THRESHOLD = $-0.1$}: scorul \texttt{score\_samples} sub $-0.1$
    declanșează o alertă \texttt{ML\_ANOMALY}; valori mai negative corespund unor puncte mai
    ușor de izolat (mai anomale);
  \item \textbf{n\_estimators = 100}: 100 de arbori de izolare, suficient pentru stabilizarea
    estimărilor conform literaturii~\cite{liu2008isolation,liu2012isolationtkdd};
  \item \textbf{reantrenare zilnică la 03:00} prin \texttt{scheduler.py} / serviciu
    \texttt{systemd}, pe ultimele 30 de zile de date.
\end{itemize}

Alegerea pragului $-0.1$ este un punct de start conservator: permite detectarea anomaliilor
moderate fără a inunda utilizatorul cu false pozitive. Calibrarea finală va depinde de datele
reale și de preferința utilizatorului față de compromisul precizie–recall.

\subsection{Mecanism de cold-start}

Deoarece modelul \gls{if} necesită minim $\sim$500 rânduri de antrenament (aprox. 14 zile la
5 secunde/citire cu pierderi de conectivitate incluse), sistemul funcționează în mod
\textbf{cold-start} în primele săptămâni, folosind reguli prag hardcodate: gaz $>$ 300,
temperatură $> 35\,°C$ sau $< 5\,°C$, mișcare în intervalul 23:00–06:00. Odată generat primul
model valid, \texttt{infer.py} comută automat pe inferența bazată pe scor.

\subsection{Rezultate ilustrative}

Figura~\figref{fig:eval} prezintă un exemplu ilustrativ al comportamentului așteptat al
detectorului de anomalii pe un interval de 24 de ore, generat cu date sintetice.

\begin{figure}[H]
  \centering
  \includegraphics[width=\singlefigure]{\pics/fig-eval-anomalii}
  \caption{Exemplu ilustrativ al scorului de anomalie pe 24\,h (date sintetice / preliminare):
    serie temporală a scorului \gls{if}, prag la $-0.1$ și puncte detectate ca anomalii (roșu).
    \textbf{Notă: datele sunt ilustrative și nu provin din operare reală de durată.}}
  \label{fig:eval}
\end{figure}

Graficul ilustrează mecanismul de detectare: scorul se menține în intervalul $[-0.05, 0.05]$
pentru comportamente normale, coboară sub $-0.1$ la injectarea unor spike-uri sintetice de
temperatură și gaz, și revine la normal odată ce perturbația dispare. Pe datele sintetice folosite
în faza de dezvoltare, parametrii actuali (\texttt{contamination=0.05}, prag $-0.1$) au produs
rezultate orientative de $F_1 \approx 0.78$ — valoare \textbf{strict ilustrativă}, obținută
pe un set mic de anomalii injectate manual, fără semnificație statistică.

% ──────────────────────────────────────────────────────────────────────
\section{Discuție: avantaje confirmate și limitări actuale}
\label{sec:discutie}
% ──────────────────────────────────────────────────────────────────────

\subsection{Avantaje confirmate ale abordării}

\textbf{Funcționare local-first.}
Verificarea funcțională a confirmat că toate fluxurile principale (monitorizare, control, alertă,
inferență \gls{ml}) funcționează integral în \gls{lan}, fără dependențe de servicii externe.
Deconectarea de la internet nu afectează niciuna dintre funcționalitățile de bază, ceea ce
aliniază sistemul cu principiile local-first~\cite{kleppmann2019localfirst} și adresează
preocupările de confidențialitate documentate în literatura \gls{gdpr}~\cite{zheng2018smarthomeprivacy,gdpr2016}.

\textbf{Cost redus al platformei hardware.}
Arhitectura bazată pe ESP32 (~3–5 EUR/unitate) și \gls{rpi} 3B+ (~35 EUR) demonstrează că
un sistem smart home cu \gls{ml} integrat este realizabil la un cost total al componentelor
sub 60 EUR, față de soluțiile comerciale care implică abonamente lunare și dependențe de cloud.

\textbf{Inferență \gls{ml} la margine (edge).}
Execuția modelului \gls{if} direct pe \gls{rpi}, cu timpi de inferență sub 5\,ms per citire,
confirmă fezabilitatea abordării edge-ML pentru scenarii \gls{iot} rezidențiale~\cite{shi2016edge,premsankar2018edgeiot}.
Modelul nesupervizat nu necesită date etichetate, avantaj esențial pentru utilizatorul rezidențial
care nu poate produce etichete manuale~\cite{chandola2009anomaly,cook2020anomalyiot}.

\textbf{Modularitate și extensibilitate.}
Separarea clară între firmware (Arduino C++), gateway (NestJS + Python) și aplicație mobilă
(React Native/Expo) facilitează adăugarea de noi noduri, senzori sau reguli de automatizare fără
modificarea celorlalte straturi.

\subsection{Limitări actuale}

\textbf{Absența datelor reale de lungă durată.}
Cea mai semnificativă limitare actuală este că modelul \gls{if} nu a fost antrenat pe date
reale provenite din operarea continuă a sistemului. Performanța sa reală — în termeni de
precizie, recall și $F_1$ — nu poate fi afirmată cu certitudine și urmează să fie măsurată
după cel puțin 30–60 de zile de colectare continuă.

\textbf{Absența \gls{tls} pe \gls{mqtt}.}
Comunicarea \gls{mqtt} curentă se realizează pe portul 1883 (plaintext), cu autentificare
simplă prin nume de utilizator și parolă. Deși acest lucru este acceptabil într-o rețea locală
controlată, adăugarea \gls{tls} (portul 8883) este planificată pentru versiunile viitoare,
în conformitate cu recomandările de securitate pentru dispozitive \gls{iot}~\cite{alrawi2019sokhome}.

\textbf{Absența testelor automate.}
Suita de teste Jest (pentru NestJS) și pytest (pentru \texttt{features.py}) nu este încă
implementată. Aceasta limitează capacitatea de regresie și face verificarea funcționalităților
dependentă de testarea manuală.

\textbf{Prototip pe breadboard.}
Asamblarea pe breadboard introduce conexiuni nesigure și variabilitate în măsurătorile analogice
(\gls{adc}), în special pentru senzorul MQ-2 și \gls{ldr}-uri. Migrarea pe un circuit imprimat
dedicat (PCB) sau cel puțin pe o placă de prototipare permanentă (perfboard) ar îmbunătăți
repetabilitatea măsurătorilor.

\textbf{Scalabilitate — nod unic interior.}
Topologia curentă include un singur nod interior. Extinderea la mai multe noduri (camere
diferite) ar necesita ajustarea modelului \gls{ml} pentru a trata fiecare locație independent
sau a construi un model agregat.

\textbf{Calibrarea senzorului MQ-2.}
Senzorul MQ-2~\cite{hanwei2020mq2} necesită o perioadă de preîncălzire (24–48 ore) și o
calibrare în aer curat (\textit{R\textsubscript{0}}) pentru a produce valori absolute de
concentrație (ppm). În prototipul curent se folosesc valori ADC brute (0–1023), suficiente
pentru detecția unui prag de alertă dar necorespunzătoare pentru cuantificarea precisă a
concentrației de gaz.

% ──────────────────────────────────────────────────────────────────────
\section{Sinteză}
\label{sec:sinteza-rezultate}
% ──────────────────────────────────────────────────────────────────────

Capitolul a prezentat metodologia de evaluare adoptată pentru prototipul de sistem smart home
local-first și a raportat starea curentă a verificărilor. Toate funcționalitățile principale
— de la publicarea datelor \gls{mqtt} și persistarea în TimescaleDB, la controlul actuatoarelor
și propagarea alertelor prin WebSocket — au fost confirmate prin testare manuală. Indicatorii
de performanță operațională (latență, consum resurse pe \gls{rpi}, timp inferență \gls{ml})
au fost prezentați ca valori ilustrative/preliminare, obținute în sesiuni punctuale de testare.

Protocolul de evaluare a detecției de anomalii cu \gls{if} a fost definit complet, inclusiv
metrica de injectare de anomalii sintetice, metricile $F_1$ și AUC-ROC și analiza sensibilității
la prag. Evaluarea cantitativă definitivă va deveni posibilă după acumularea datelor de operare
reală. Limitările principale identificate — absența \gls{tls} pe \gls{mqtt}, lipsa testelor
automate și dependența de date de durată pentru model — constituie direcții concrete de
dezvoltare ulterioară, discutate în \chapref{ch:concluzii}.
